
\section{Evaluation Agenda}


\begin{figure}[H]
\centering
\begin{tikzpicture}[every node/.style={figlabel,align=center}, cell/.style={figbox, minimum width=2.85cm, minimum height=.95cm}]
\path[figpanel] (-1.1,-0.95) rectangle (11.0,2.55);
\node[figtitle, anchor=west] at (-0.65,2.08) {EVALUATION};
\node[cell, fill=BitGray] at (0.4,1.15) {Baseline};
\node[cell, fill=BitBlue!8, draw=BitBlue!65] at (3.45,1.15) {Manifest only};
\node[cell, fill=BitTeal!8, draw=BitTeal!70] at (6.5,1.15) {Minimal core};
\node[cell, fill=BitRose!10, draw=BitRose!62] at (9.55,1.15) {Bloated / stale};
\node[cell] at (0.4,-0.1) {success\\cost};
\node[cell] at (3.45,-0.1) {instruction\\fit};
\node[cell] at (6.5,-0.1) {restart\\drift};
\node[cell] at (9.55,-0.1) {failure\\modes};
\end{tikzpicture}
\caption{The evaluation should compare useful state, no state, and harmful state.}
\label{fig:evaluation-matrix}
\end{figure}


The right evaluation target is not "does any context file help?" The right target is narrower: under what conditions does a compact, validated state layer improve repository work, and when does it become overhead?

\subsection{Conditions}

The core conditions are baseline (no repository state), manifest-only, minimal Continuity, minimal Continuity with hooks, full topology, and failure modes such as stale state, bloated context, and conflicting local instructions.

\subsection{Benchmarks and tasks}

Existing repository benchmarks such as \SWEbench{} remain useful for patch correctness \cite{Jimenez2024SWEbench}. They do not, however, measure resumption, repeated-error avoidance, or multi-agent collision. CONTINUITY therefore needs a custom task family with interruptions, stale-state traps, cross-branch handoffs, and malicious or noisy context.

\subsection{Metrics}


\begin{table}[H]
\centering
\caption{Primary metric families for a CONTINUITY evaluation.}
\label{tab:evaluation-metrics}
\small
\rowcolors{2}{BitGray}{white}
\begin{tabularx}{\linewidth}{>{\raggedright\arraybackslash}p{2.9cm}X}
\rowcolor{BitNavy!8}
\textbf{Metric family} & \textbf{Examples} \\
Task outcome & success rate, pass@1, regression count, patch correctness \\
Efficiency & input tokens, output tokens, tool calls, wall-clock time, cost \\
Continuity & resumption latency, repeated-error rate, stale-state incidence, reconciliation completeness \\
Multi-agent & claim conflict rate, duplicate work, merge conflicts, handoff success \\
Governance & instruction compliance, review time, smell prevalence, override frequency \\
Security & prompt-injection obedience, secret leakage, unauthorized action attempts \\
\end{tabularx}
\end{table}


The most interesting metrics are often not the headline ones. Task success matters, but so do resumption latency, repeated-error rate, evidence completeness, human review time, and collision frequency.

\subsection{Ablations}

Ablations should remove \statefile{STATE.md}, \statefile{TODO.md}, \statefile{LEARNINGS.md}, scoped loading, evidence requirements, and validation hooks. Additional ablations should intentionally inject context bloat, stale facts, and local instruction conflicts.

\subsection{Expected shape of the results}

The expected positive result is modest: faster restarts, better handoff, cleaner audit trails, and fewer repeated failures. The expected negative result is equally important: larger context topologies should cost more and may perform worse. A credible paper should publish both.
